% Abstände – bank/abstaende/ (zusammenbau v0.4, Heft)
\einheitenkopf[t7]{Abstände}
\setcounter{aufgabe}{64}

% Anlauf (ohne Prüfkennung)
\begin{aufgabe}{Abstand zweier Punkte – Anlauf}
\begin{teile}
\teil Abstand wovon zu was? Wie weit ist die Mastspitze S(3 | 2 | 7) vom Punkt T(0 | 6 | 7) entfernt? Kreuze das Verfahren an. \\ \kreuz{Punkt–Punkt: Betrag des Verbindungsvektors} \\ \kreuz{Punkt–Ebene: Hessesche Normalform} \\ \kreuz{Punkt–Gerade: Lotfußpunkt}
\teil A(2 | 1 | 3), B(4 | 3 | 4) – wie weit voneinander entfernt? d = \leerfeld
\teil Ein Segelflugzeug gleitet geradlinig in 40 s von A(200 | 300 | 1\,000) nach B(800 | 1\,500 | 600); 1 LE = 1 m. Wie schnell ist es in km/h? v = \leerfeld km/h
\end{teile}
\end{aufgabe}

\begin{aufgabe}{Abstand zweier Punkte – Prüfungsaufgaben}
\begin{teile}
\teil Eine Drohne fliegt geradlinig in 4 s von P(5 | 20 | 30) nach Q(65 | 100 | 30); 1 LE = 1 m. Wie schnell ist sie in km/h? \hfill (Abitur 2018 GK) \\ v = \leerfeld km/h
\teil Ein Klettergerüst hat die Kanten AB mit A(0 | 1 | 4), B(4 | 1 | 4) und CD mit C(2 | 7 | 1), D(6 | 7 | 1); 1 LE = 1 m. Ein Seil zwischen den Mittelpunkten der beiden Kanten ist 30 \% länger als deren Abstand. Wie lang ist das Seil? \hfill (Abitur 2018 GK) \\ \leerfeld[m]
\teil Eine rechteckige Terrasse hat die Ecken A(2 | 1 | 0), B(9 | 1 | 0), C(9 | 7 | 0) und D(2 | 7 | 0); 1 LE = 1 m. Ein Gartenschlauch ist am Wasserhahn P(3 | 2 | 2) befestigt und 9 m lang. Erreicht er jede Stelle der Terrasse? \hfill (Abitur 2023 GK) \\ \janein
\teil Das Dreieck ABC mit A(0 | 0 | 0), B(2 | 1 | 2) und C(t | −2t | 0), t positiv, ist rechtwinklig in A und gleichschenklig. Welcher Wert von t? \hfill (Abitur 2026 GK) \\ t = \leerfeld
\teil Gegeben sind P(1 | 3 | −2) und Q(9 | y | 13) mit unbekannter Koordinate y. Weise nach, dass P und Q mindestens den Abstand 17 haben. \hfill (Abitur 2019 GK)
\teil Gegeben sind R(−1 | 4 | 2) und S(4 | −8 | z) mit unbekannter Koordinate z. Weise nach, dass R und S mindestens den Abstand 13 haben. \hfill (Abitur 2019 GK)
\end{teile}
\end{aufgabe}

% Anlauf (ohne Prüfkennung)
\begin{aufgabe}{Abstand Punkt–Ebene – Anlauf}
\begin{teile}
\teil Hin oder zurück? Wie weit ist P(2 | −1 | 3) von der Ebene E: $2x - 2y + z = 1$ entfernt? Kreuze an, rechne nicht. \\ \kreuz{Abstand berechnen} \\ \kreuz{Punkt zum vorgegebenen Abstand bestimmen}
\teil Abstand von P(4 | 0 | 3) zur Ebene E: $2x - y + 2z = 5$? d = \leerfeld
\teil Lotgerade durch P(−1 | 0 | 5) zur Ebene E: $3x - 4z = 1$? $\vec x = $ \leerfeld
\end{teile}
\end{aufgabe}

\begin{aufgabe}{Abstand Punkt–Ebene – Prüfungsaufgaben}
\begin{teile}
\teil Die Ebene E: $3y + 4z = 3$ enthält den Punkt (2 | 1 | 0). Gesucht ist ein Punkt P, dessen Spiegelpunkt an E von P den Abstand 14 hat. Gib einen solchen Punkt an. \hfill (Abitur 2025 GK) \\ P( \leerfeld | \leerfeld | \leerfeld )
\end{teile}
\end{aufgabe}

% Anlauf (ohne Prüfkennung)
\begin{aufgabe}{Lotfußpunkt – Anlauf}
\begin{teile}
\teil Was sagt die Gleichung $\overrightarrow{PF} \circ \vec r = 0$ in einem Lösungsweg? Kreuze an. \\ \kreuz{Punkt auf der Geraden} \\ \kreuz{Lotbedingung (Skalarprodukt null)} \\ \kreuz{Abstandsgleichheit}
\teil Lotfußpunkt F von P(4 | 2 | 3) auf g: $\vec x = (2 | 1 | 0) + t \cdot (2 | -1 | 2)$ – und der Abstand von P zu g? F( \leerfeld | \leerfeld | \leerfeld ), d ≈ \leerfeld
\teil Die Gleichungen I $\overrightarrow{OX} = \overrightarrow{OM} + s \cdot \vec u$ und II $\overrightarrow{PX} \circ \vec u = 0$ liefern gemeinsam einen Wert von s. Welche Bedeutung haben X und die Länge $|\overrightarrow{PX}|$?
\end{teile}
\end{aufgabe}

\begin{aufgabe}{Lotfußpunkt – Prüfungsaufgaben}
\begin{teile}
\teil An einem Zeltgestänge liefern die Gleichungen I $\overrightarrow{OQ} = \overrightarrow{OA} + r \cdot \overrightarrow{AB}$ und II $\overrightarrow{EQ} \circ \overrightarrow{AB} = 0$ gemeinsam einen Wert von r. Welche Bedeutung hat der Punkt Q für diesen Wert? \hfill (Abitur 2026 GK)
\teil Bei einem Carport soll von der Mitte M des Balkens AB eine möglichst kurze Strebe zum Balken CD führen; sie trifft CD im Punkt R. Beschreibe, wie man die Koordinaten von R ermitteln kann. \hfill (Abitur 2022 GK)
\teil Die Kante AB eines Körpers liegt in der xy-Ebene und in der Ebene L: $2x + y + z = 8$. Zum Punkt Q(0 | 3 | 0) ist vorgelegt: I g: $\vec x = (0 | 3 | 0) + t \cdot (2 | 1 | 0)$; II $2 \cdot 2t + 3 + t = 8 \Rightarrow t = 1 \Rightarrow R(2 | 4 | 0)$; III $\sqrt{(0 - 2)^2 + (3 - 4)^2} \approx 2{,}24$. Erläutere die Schritte und gib die geometrische Bedeutung des Ergebnisses an. \hfill (Abitur 2023 GK)
\teil Pyramide über dem Quadrat ABCD mit A(1 | 0 | 0), B(5 | 0 | 0), C(5 | 4 | 0), D(1 | 4 | 0) und der Spitze E(1 | 4 | 2); B, D, E liegen in einer Symmetrieebene. Vorgelegt ist ein Rechenweg für die kürzeste Linie von A über die Kante BE nach C: $P(5 - 4t | 4t | 2t)$; $\overrightarrow{PA} \circ \overrightarrow{PB} = 0 \Leftrightarrow t = \tfrac{4}{9}$; $2 \cdot |\overrightarrow{PA}| \approx 5{,}96$. Erläutere das Vorgehen. \hfill (Abitur 2021 GK)
\teil Ein Mähroboter mit kreisförmiger Unterseite (Radius 0,3 m) fährt mit seinem Mittelpunkt auf der Geraden g mit Richtungsvektor (3 | 4 | 0) auf den Rand k der Rasenfläche zu; g trifft k im Punkt Q(6 | 8 | 0) unter dem Winkel 30°; 1 LE = 1 m. Der Roboter wendet, sobald seine Unterseite k berührt. Wo ist dann sein Mittelpunkt S? Nutze eine Skizze. \hfill (Abitur 2021 GK) \\ S( \leerfeld | \leerfeld | \leerfeld )
\end{teile}
\end{aufgabe}

% Anlauf (ohne Prüfkennung)
\begin{aufgabe}{Abstand als Argument – Anlauf}
\begin{teile}
\teil Senkrecht oder schräg? Die Strecke verbindet P(1 | 2 | 4) mit dem Punkt Q(1 | 2 | 0) der xy-Ebene. Kreuze an, rechne nicht. \\ \kreuz{senkrecht: die Strecke ist der Abstand} \\ \kreuz{schräg: die Strecke ist länger als der Abstand}
\teil Die zwei Scheiben eines Doppelfensters liegen in den parallelen Ebenen $E_1$: $2x - y + 2z = 1$ und $E_2$: $2x - y + 2z = 7$; P(0 | −1 | 0) liegt in $E_1$, Q(1 | 1 | 3) in $E_2$. Ist PQ kleiner, größer oder genauso groß wie der Abstand der Scheiben? Begründe.
\teil Die Abbildung zeigt eine Rampe im Querschnitt als Strecke GJ auf der Geraden t und eine Lampe F. F hat von t etwa den Abstand 3,5. Begründe mit einer Zeichnung, dass kein Punkt der Rampe von F den Abstand 3,5 hat.

\begin{ksys}[xmin=0,xmax=8,ymin=0,ymax=8]
\gerade{-1}{8}{t}
\punkt{4}{4}{G}
\punkt{7}{1}{J}
\punkt{0}{3}{F}
\end{ksys}
\end{teile}
\end{aufgabe}

\begin{aufgabe}{Abstand als Argument – Prüfungsaufgaben}
\begin{teile}
\teil B(1 | 0 | 4) liegt in der Ebene $L_1$: $x + 2y + 2z = 9$, E(1 | 0 | 1) in der parallelen Ebene $L_2$: $x + 2y + 2z = 3$. Ist der Abstand von B und E kleiner, größer oder genauso groß wie der Abstand der Ebenen? Begründe. \hfill (Abitur 2023 GK)
\teil Ein Dachfirstpunkt F ist 3 m vom Punkt G einer schrägen Glasfläche entfernt; die Strecke FG steht nicht senkrecht auf der Glasfläche. Begründe, dass der Abstand von F zur Ebene der Glasfläche kleiner als 3 m ist. \hfill (Abitur 2023 GK)
\teil Eine Seilrutsche führt vom Startpunkt $S_1$ zum tiefer gelegenen Zielpunkt $S_2$. Auf der Karte erscheint als Länge der Anlage der horizontale Abstand von $S_1$ und $S_2$. Begründe ohne Rechnung, dass diese Länge kürzer ist als die Seilstrecke $S_1S_2$. \hfill (Abitur 2021 GK)
\teil Gegeben sind P(2 | 1 | 0), Q(8 | 1 | 0) und R(2 | 5 | 0). Gib zwei Punkte an, die von P, Q und R gleich weit entfernt sind. \hfill (Abitur 2019 GK) \\ ( \leerfeld | \leerfeld | \leerfeld ), ( \leerfeld | \leerfeld | \leerfeld )
\teil Ein Zeltdach hat die Spitze S(2 | 2 | 9); F(3 | 2 | 6) und G(2 | 3 | 6) liegen auf zwei Dachkanten. Die Geraden SF und SG schneiden die xy-Ebene in $Q_1$ und $Q_2$. In welchem Verhältnis steht der Abstand von $Q_1$ und $Q_2$ zum Abstand von F und G? Begründe, ohne $Q_1$ und $Q_2$ zu berechnen. \hfill (Abitur 2022 GK)
\teil Eine Pyramide hat die Spitze S(1 | 1 | 8); F(2 | 1 | 6) und G(1 | 2 | 6) liegen auf zwei Seitenkanten. Die Geraden SF und SG schneiden die xy-Ebene in $Q_1$ und $Q_2$. In welchem Verhältnis steht der Abstand von $Q_1$ und $Q_2$ zum Abstand von F und G? Begründe, ohne $Q_1$ und $Q_2$ zu berechnen. \hfill (Abitur 2022 GK)
\end{teile}
\end{aufgabe}
